% Joint PDF surface for Example 2 (water levels) in M3f-Joint-Distributions.tex.
% Precomputed here because pgfplots surfaces are slow to render on every lecture build.
% Page 1: surface only; page 2: adds the slice at x = 0.5 used in part (b).
% Rebuild with:
%   ./build-plots.sh water-levels-joint-pdf
\documentclass[smaller]{beamer}
\usepackage{tikz}
\usepackage{pgfplots}
\pgfplotsset{compat=newest}

\setbeamertemplate{navigation symbols}{}
\setbeamertemplate{footline}{}
\setbeamertemplate{headline}{}

\begin{document}
\begin{frame}[plain]
  \centering
  \pgfplotsset{
    colormap={whitered}{color(0cm)=(white); color(1cm)=(purple!75!blue)}
  }

  \begin{tikzpicture}[scale=.6,
    declare function={f(\x,\y)=6/5*(\x + \y^2);}]
    \begin{axis}[
      colormap name=whitered,
      width=14cm,
      view={-30}{30},
      enlargelimits=false,
      grid=major,
      domain=0:1,
      y domain=0:1,
      zmin=0, zmax=2.4,
      samples=21,
      xlabel=$x$,
      ylabel=$y$,
      zlabel={$f(x,y)$},
      colorbar,
      colorbar style={
        at={(1.1,0)},
        anchor=south west,
        height=0.25*\pgfkeysvalueof{/pgfplots/parent axis height},
        title={$f(x,y)$}
      }
      ]
      \addplot3 [surf, faceted color=black!30] {f(x,y)};

      % Slice at x = 0.5 (part b): area under f(0.5, y) and the curve itself
      \only<2->{
        \addplot3 [draw=none, fill=red, fill opacity=.2, samples=21, samples y=0, domain=0:1]
          (0.5, x, {f(0.5,x)}) -- (axis cs:0.5,1,0) -- (axis cs:0.5,0,0) -- cycle;
        \addplot3 [red, very thick, samples=21, samples y=0, domain=0:1] (0.5, x, {f(0.5,x)});
        \node[red] at (axis cs:0.5,1,{f(0.5,1)}) [pin={[red]90:$f(0.5,y)$}] {};
      }
    \end{axis}
  \end{tikzpicture}
\end{frame}
\end{document}
